\section{Related Work}\label{sec:related_work}

\subsection{Physics-based motion imitation}

DeepMimic demonstrated reference-guided reinforcement learning for simulated
character skills~\cite{deepmimic}. Perpetual Humanoid Control developed robust
humanoid imitation for real-time simulated avatars~\cite{phc}. These approaches
motivate learning a feedback controller rather than directly replaying a
kinematic trajectory. PULSE learns a motion representation for humanoid control, while OmniH2O,
ExBody2, and GMT study large-scale control on robot platforms~\cite{pulse,omnih2o,exbody2,gmt}.
ASE learns reusable latent skills and downstream task policies rather than
requiring exact clip tracking~\cite{ase}. These objectives provide context,
but their published scores do not measure our clip/body replay task.
We implement our experiments in ProtoMotions with Isaac Gym and PPO
optimization~\cite{protomotions,isaacgym,ppo}.

\subsection{Morphology-conditioned control and retargeting}

Body-shape variation has already been studied in physics-based control:
Won and Lee learned a parametric controller that accommodates variation in
character height, weight, and proportions~\cite{won2019}. A separate line of
work learns policies that transfer across changes in agent \emph{topology} --
limb count and skeletal structure -- through shared modular
sub-policies~\cite{smp2020}, which is a different axis from the continuous
anthropometric variation studied here. Morphology conditioning is therefore
an established idea, and our work should not be interpreted as introducing it.
We instead examine one controller over a large set of paired motions and
bodies, with shape identity maintained throughout data generation,
simulation, and evaluation.

Kinematic retargeting and dynamic tracking address different requirements.
SMPL provides a common parameterized body representation~\cite{smpl}, and
SMPLSim constructs simulation-ready assets from these
models~\cite{smplsim}. A shared skeleton makes poses comparable, but does not
make global positions or required actuation independent of shape. Matching a
reference to its intended simulation asset is consequently part of the task
definition, not an optional evaluation convenience.

\subsection{Text-motion data, shape-conditioned generation, and refinement}

AMASS consolidates motion capture in a common body-model
representation~\cite{amass}. HumanML3D associates processed motion segments
with natural-language descriptions~\cite{humanml3d}. We inherit those source
identities and descriptions rather than claiming new motion capture or
shape-specific language annotation.

HUMOS generates identity-conditioned human motion using a cycle-consistent
learning formulation~\cite{humos}. We use its pretrained model to obtain
shape-specific references, not as the physics controller and not as a
text-conditioned policy. The resulting sequences still require conversion to
the simulator's skeleton and coordinate conventions, and their feasibility
must be assessed in that simulator. Our refinement is a deliberately limited
kinematic post-process; it is not a new generative model or a full dynamic
trajectory optimizer.

The controller uses standard Transformer operations~\cite{attention}.
Learned slot and type embeddings identify temporal positions and observation
sources. An ablation also investigates actor-only morphology modulation with
adaLN-Zero, adapting the zero-initialized conditioning design used in
Diffusion Transformers~\cite{dit}. That experiment tests a conditioning
mechanism, not diffusion-based action generation.

\subsection{Headline capability and evidence comparison}
Table~\ref{tab:capabilities} separates published support from demonstrated
evaluation coverage. In particular, PHC's Appendix B describes shape-varying
humanoids and training with sampled AMASS shapes; its quantitative comparisons
use the mean neutral SMPL body~\cite{phc}. It is therefore inaccurate to
classify PHC as incapable of body variation. Won and Lee directly establish
parametric body control~\cite{won2019}. Our distinction is the paired corpus
and systematic inside/outside replay, not priority for morphology conditioning.

\begin{table}[p]
\centering\footnotesize
\caption{Published capabilities and evaluation scope, not a performance
ranking. ``Not reported'' concerns the cited source and the specified
unseen-body test, not impossibility. No external system has been rerun on our
matched references, assets, horizon, and aggregation. Software documentation
for ProtoMotions 3 is distinguished from paper evaluations and from the
implementation used in this study.}\label{tab:capabilities}
\begin{tabularx}{\linewidth}{@{}p{0.13\linewidth}YYY@{}}
\toprule
System / source & Task and motion evidence & Body capability & Evaluation relevance \\
\midrule
PHC~\cite{phc} & Reference tracking and recovery; 11,313 train and 140 test AMASS sequences. & SMPL shape/gender variation; sampled shapes and body-normalized references described. & Quantitative motion tests use mean neutral body; shape variation shown qualitatively. No systematic inside/outside unseen-body panel reported. \\[0.5em]
ASE~\cite{ase} & Latent skill pretraining and downstream tasks; 187 clips. Exact reference tracking is not its objective. & Published experiments use a sword/shield humanoid; a shared shape-conditioned body sweep is not reported. & Task transfer, skill coverage, and recovery evaluated; unseen-body clip-tracking test not reported. \\[0.5em]
ProtoMotions 3~\cite{protomotions} & Framework documents large-scale tracking, distributed training, retargeting, and generative policies. & Supports different configured characters/robots; this alone does not demonstrate one policy across shape variants. & Official README is a software capability source, not a matched multi-body benchmark. Our simulator/PPO/tracking infrastructure is inherited. \\[0.5em]
Won \& Lee~\cite{won2019} & Parametric physical character control with adaptive shape sampling. & One parametric controller accommodates height, weight, and proportion changes; new and changing bodies demonstrated. & Direct prior body-control capability; experiments differ from our large clip-by-body imitation panels. \\[0.5em]
HUMOS~\cite{humos} & Identity-conditioned motion generation with physical-consistency objectives. & Shape-conditioned motions; supplies our pretrained reference generator. & Not an online feedback tracker; its motion-generation scores are not tracking success under our simulator protocol. \\[0.5em]
This study & Shared tracking; 18,855 training identities, 1,048 held-out test identities. & 128 trained bodies; references and assets share identity. & Test: one trained body per clip. Unseen: 233 known clips $\times$ each of 128 inside and 128 outside bodies; all shapes retained. Conditioning mechanism remains unisolated. \\
\bottomrule
\end{tabularx}
\end{table}

\paragraph{Inherited functions and study additions.}
ProtoMotions provides the simulation abstraction, PPO training, motion-library
and imitation infrastructure. HUMOS supplies the pretrained generator;
SMPL/SMPLSim supply the body parameterization and asset construction. Our
study adds the paired corpus assembly and identity checks, reference
processing and fidelity audit, the evaluated temporal slot/type controller
configuration, streamed clip residency, and exposure-separated evaluation.
These are system and empirical contributions built on existing methods,
not claims to invent attention, PPO, body models, or motion imitation.
A directly comparable performance claim requires a matched external replay;
none is made here.
